\section{Implementierung der CI/CD Pipeline}

\subsection{Erweiterung der CI/CD-Pipeline zur Ausführung automatisierter API-Tests}

Zur Sicherstellung der Funktionsfähigkeit des gesamten Systems wurden automatisierte API-Tests mithilfe von Playwright in die bestehende \texttt{GitHub Actions}-Pipeline integriert. Ziel ist es, die korrekte Verarbeitung von SOAP-Anfragen wie \texttt{addUser}, \texttt{updateUser}, \texttt{deleteUser} und \texttt{getUserStatus} automatisiert zu prüfen.

\paragraph{Initialisierungsskript}

Um in der Pipeline eine konsistente Testumgebung sicherzustellen, wurde ein separates Initialisierungsskript \texttt{init-ldap.sh} erstellt. Dieses Skript stellt sicher, dass der OpenLDAP-Container vollständig gestartet ist und initialisiert anschliessend die Datenbank mit einer vordefinierten Struktur sowie Testdaten. Dazu gehören die benutzerdefinierte Schema-Datei, die organisatorische Struktur, Gruppen und Beispielbenutzer.

\begin{lstlisting}[language=bash, caption={Von ChatGPT generiertes Initialisierungsskript für LDAP-Daten in der CI/CD-Pipeline}]
#!/bin/bash

set -e
set -o pipefail

echo ">>> Waiting for OpenLDAP container to be ready..."

timeout=30
elapsed=0

until docker exec openldap ldapsearch -Y EXTERNAL -H ldapi:/// -b "" -s base "(objectclass=*)" > /dev/null 2>&1; do
  sleep 1
  elapsed=$((elapsed+1))
  if [ "$elapsed" -ge "$timeout" ]; then
    echo " OpenLDAP did not become ready within ${timeout} seconds."
    exit 1
  fi
done

echo "OpenLDAP is ready."

docker exec openldap ls -la /container/service/slapd/assets
docker exec openldap ldapadd -Y EXTERNAL -H ldapi:/// -f /container/service/slapd/assets/schema.ldif
docker exec openldap ldapadd -Y EXTERNAL -H ldapi:/// -f /container/service/slapd/assets/01_ou_structure.ldif
docker exec openldap ldapadd -Y EXTERNAL -H ldapi:/// -f /container/service/slapd/assets/02_groups.ldif
docker exec openldap ldapadd -Y EXTERNAL -H ldapi:/// -f /container/service/slapd/assets/03_users.ldif

echo " LDAP setup completed successfully."
\end{lstlisting}

\paragraph{Erklärung der Pipeline-Schritte}

Der neue Pipeline-Job \texttt{api-tests} wird nach dem erfolgreichen Build gestartet. Die einzelnen Schritte werden im Folgenden erläutert:

\begin{enumerate}
    \item \textbf{Code Checkout:} Der Quellcode des Projekts wird aus dem Repository geladen.
    \item \textbf{Debug-Ausgabe (optional):} Ausgabe des aktuellen Verzeichnisinhalts zur Fehlersuche, z.\,B. falls Pfade nicht stimmen.
    \item \textbf{Docker Compose Installation:} Da GitHub Actions kein Docker Compose vorinstalliert hat, wird dieses Werkzeug manuell installiert. Es wird später für das Starten von Applikation und LDAP verwendet.
    \item \textbf{Node.js Setup:} Die für Playwright benötigte Node.js-Version (hier: Version 20) wird bereitgestellt.
    \item \textbf{Installation der JavaScript-Abhängigkeiten:} Mit \texttt{npm ci} werden alle im Projekt definierten JavaScript-Module installiert, die für die Ausführung der Tests benötigt werden.
    \item \textbf{JDK Setup:} Da die Applikation in Java entwickelt ist, wird das Java Development Kit in Version 17 eingerichtet.
    \item \textbf{Build der Applikation:} Das Projekt wird mit \texttt{mvnw clean install} gebaut. Die Tests werden hierbei übersprungen, da sie separat mit Playwright abgedeckt werden.
    \item \textbf{Start der benötigten Container:} Die Applikation sowie der LDAP-Server werden mittels \texttt{docker-compose up -d} gestartet.
    \item \textbf{Verfügbarkeit prüfen:} Es wird mit einem einfachen \texttt{curl}-Befehl regelmässig geprüft, ob die Applikation über ihren Heartbeat-Endpunkt erreichbar ist.
    \item \textbf{LDAP mit Testdaten befüllen:} Das oben beschriebene Initialisierungsskript wird ausgeführt, um das LDAP-System für die Tests vorzubereiten.
    \item \textbf{Ausführen der Playwright-Tests:} Die eigentlichen API-Tests werden mit dem Befehl \texttt{npm run test:api} ausgeführt.
    \item \textbf{HTML Report:} Die ausgeführten API-Tests werden als HTML-Report gespeichert.
    \item \textbf{Aufräumen:} Unabhängig vom Testergebnis werden alle gestarteten Container am Ende wieder heruntergefahren.
\end{enumerate}

\begin{lstlisting}[language=yaml, caption={GitHub Actions Job zur Ausführung der API-Tests}]
  api-tests:
    needs: build
    runs-on: ubuntu-latest
    steps:
      - name: Check out code
        uses: actions/checkout@v4

      - name: Show current directory structure (debug)
        run: |
          pwd
          ls -l

      - name: Install Docker Compose
        run: |
          sudo curl -L "https://github.com/docker/compose/releases/latest/download/docker-compose-$(uname -s)-$(uname -m)" -o /usr/local/bin/docker-compose
          sudo chmod +x /usr/local/bin/docker-compose
          docker-compose version

      - name: Set up Node.js for Playwright
        uses: actions/setup-node@v4
        with:
          node-version: '20'

      - name: Install Node.js dependencies
        run: npm ci
        working-directory: ${{ github.workspace }}

      - name: Set up JDK
        uses: actions/setup-java@v4
        with:
          distribution: 'temurin'
          java-version: '17'

      - name: Build project
        run: chmod +x ./mvnw && ./mvnw clean install -DskipTests
        working-directory: ${{ github.workspace }}

      - name: Start Docker services
        run: docker-compose up -d
        working-directory: ${{ github.workspace }}

      - name: Wait for app and LDAP to be ready
        run: |
          until curl -s -o /dev/null -w "%{http_code}" http://localhost:8080/api/heartbeat | grep -q "200"; do
            echo "Waiting for app and LDAP to be ready..."
            sleep 2
          done
          echo "App and LDAP are ready!"

      - name: Initialize LDAP with test data
        run: bash ./scripts/init-ldap.sh
        working-directory: ${{ github.workspace }}

      - name: Run Playwright API tests
        run: npm run test:api
        working-directory: ${{ github.workspace }}

      - name: Upload Playwright Report
        if: always()
        uses: actions/upload-artifact@v4
        with:
          name: playwright-report
          path: playwright-report

      - name: Shutdown and cleanup
        if: always()
        run: docker-compose down -v
        working-directory: ${{ github.workspace }}
\end{lstlisting}

\subsubsection{Automatisierte Ausführung durch Cronjob}
Damit die Tests auch unabhängig von einem \texttt{push} oder \texttt{commit} ausgeführt werden, wurde die Pipeline um einen zeitbasierten Auslöser (\textit{Cronjob}) erweitert. Dadurch wird der komplette Workflow automatisch zweimal täglich gestartet – einmal morgens um 08:00 Uhr und einmal nachmittags um 16:00 Uhr (MEZ). Die Konfiguration befindet sich im oberen Bereich der Pipeline und sieht wie folgt aus:

\begin{lstlisting}[language=yaml, caption={Cronjob-Integration in der GitHub Actions Pipeline}]
on:
  push:
    branches:
      - main
  schedule:
    - cron: '0 6 * * *'
    - cron: '0 14 * * *'
  workflow_dispatch:
\end{lstlisting}

\subsubsection{Problem während der Entwicklung: Synchronisationsfehler zwischen Applikation und LDAP}

Während der Entwicklung kam es wiederholt zu einem Problem in der CI/CD-Pipeline: Die automatisierten Tests wurden bereits gestartet, bevor die Applikation vollständig hochgefahren war oder bevor eine stabile Verbindung zwischen Applikation und dem OpenLDAP-Container bestand. Dies führte zu instabilen Testergebnissen oder Verbindungsfehlern.

\paragraph{Lösung: Heartbeat-Endpoint mit LDAP-Prüfung}

Zur Lösung dieses Problems wurde ein dedizierter HTTP-Endpunkt \texttt{/api/heartbeat} in der Applikation implementiert. Dieser überprüft nicht nur, ob der Webserver läuft, sondern stellt explizit sicher, dass auch eine Verbindung zum LDAP-Server erfolgreich aufgebaut werden kann. Erst wenn diese Prüfung erfolgreich ist, startet die Pipeline mit den automatisierten Tests.

Die Implementierung verwendet den \texttt{LdapTemplate} von Spring, um eine einfache LDAP-Suchanfrage durchzuführen. Bei Erfolg antwortet der Endpunkt mit \texttt{HTTP 200 OK}, andernfalls mit \texttt{503 Service Unavailable}.

\begin{lstlisting}[language=Java, caption={HeartbeatController zur Verfügbarkeitsprüfung von Applikation und LDAP}]
package ch.pax.ecohub.endpoints;

import lombok.RequiredArgsConstructor;
import org.springframework.http.ResponseEntity;
import org.springframework.ldap.core.AttributesMapper;
import org.springframework.ldap.core.LdapTemplate;
import org.springframework.web.bind.annotation.GetMapping;
import org.springframework.web.bind.annotation.RestController;

@RestController
@RequiredArgsConstructor
public class HeartbeatController {

    private final LdapTemplate ldapTemplate;

    @GetMapping("/api/heartbeat")
    public ResponseEntity<String> getHeartbeat() {
        try {
            ldapTemplate.search(
                "",
                "(objectClass=*)",
                (AttributesMapper<Object>) attributes -> null
            );
            return ResponseEntity.ok("OK");
        } catch (Exception e) {
            return ResponseEntity.status(503).body("LDAP not reachable: " + e.getMessage());
        }
    }
}
\end{lstlisting}

\paragraph{Verwendung in der Pipeline}

In der \texttt{pipeline.yml} wurde ein Warte-Mechanismus eingebaut, der in regelmässigen Abständen die Erreichbarkeit dieses Heartbeat-Endpunkts prüft. Erst wenn eine gültige \texttt{HTTP 200}-Antwort empfangen wird, wird mit der Initialisierung des LDAP und dem Ausführen der Tests fortgefahren. Dadurch wird sichergestellt, dass die Tests nur unter stabilen Bedingungen starten.

\begin{lstlisting}[language=yaml, caption={Warten auf Heartbeat-Endpunkt in der Pipeline}]
- name: Wait for app and LDAP to be ready
  run: |
    until curl -s -o /dev/null -w "%{http_code}" http://localhost:8080/api/heartbeat | grep -q "200"; do
      echo "Waiting for app and LDAP to be ready..."
      sleep 2
    done
    echo "App and LDAP are ready!"
\end{lstlisting}
